\chapter{Excessive Sweating}

\section*{Definition and Overview}
Excessive sweating, known medically as hyperhidrosis, is sweating that is far greater than needed to regulate body temperature. It can affect the whole body or, more commonly, specific areas such as the armpits, hands, feet, and face. Primary hyperhidrosis usually begins in childhood or adolescence, affects specific areas, and has no underlying medical cause, while secondary hyperhidrosis is caused by another condition or medicine and tends to affect the whole body. Excessive sweating is treatable, and no one needs to simply live with it.

\section*{Epidemiology}
Hyperhidrosis affects about 1 to 3 per cent of the population, making it more common than is often assumed. It begins in childhood or adolescence in most cases and often runs in families. It affects both sexes and can cause significant embarrassment, avoidance of social situations, and reduced quality of life. Many people do not seek help, believing nothing can be done.

\section*{Aetiology and Pathophysiology}
Sweating is controlled by the sympathetic nervous system, and hyperhidrosis reflects overactivity of this system.
\begin{itemize}
    \item \textbf{Primary hyperhidrosis:} overactive sweat glands, usually in the armpits, palms, soles, and face, with no underlying disease; there is a genetic component
    \item \textbf{Secondary hyperhidrosis:} caused by an underlying condition or medicine, including an overactive thyroid, diabetes, infection, menopause, anxiety, and some medicines
    \item \textbf{Mechanism:} the sympathetic nerves stimulate the sweat glands excessively, sometimes in response to trivial triggers or with no trigger at all
    \item \textbf{Triggers:} heat, exercise, stress, and spicy food worsen sweating
    \item \textbf{Effects:} the constant wetness irritates the skin and can lead to infection, and the social effects are significant
\end{itemize}

\section*{Clinical Features}
The pattern of sweating helps distinguish the type.
\begin{itemize}
    \item Visible, excessive sweating of the armpits, hands, feet, or face
    \item Sweating that occurs without heat or exertion
    \item Sweating that soaks clothing or interferes with activities
    \item Onset in childhood or adolescence with a family history suggests primary hyperhidrosis
    \item Whole-body sweating, especially at night, suggests a secondary cause
    \item Skin changes: soft, peeling skin on the feet, and increased risk of fungal and bacterial infections
\end{itemize}

\section*{Differential Diagnosis}
Other conditions cause or worsen sweating.
\begin{itemize}
    \item Overactive thyroid, with weight loss, palpitations, and heat intolerance
    \item Menopause, with hot flushes and night sweats
    \item Diabetes, including low blood sugar
    \item Infections, including tuberculosis, with night sweats
    \item Anxiety and panic disorders
    \item Cancers, including lymphoma, with drenching night sweats
    \item Medicines, including antidepressants and some blood pressure drugs
\end{itemize}

\section*{When to Seek Medical Care}
People bothered by excessive sweating should see a GP, both for treatment and to exclude a secondary cause. Medical care is needed for sudden, severe, whole-body sweating, and for night sweats with weight loss or fever.

\textbf{Call 000 (or local emergency services) for:}
\begin{itemize}
    \item Sweating with chest pain, which may indicate a heart attack
    \item Cold, clammy skin with collapse or faintness
    \item Sweating with confusion or very low blood sugar
    \item High fever with sweating in a very unwell person
\end{itemize}

\section*{Diagnosis and Investigations}
Diagnosis is clinical, with tests to exclude secondary causes.
\begin{itemize}
    \item \textbf{History:} the pattern, triggers, and timing of sweating, and any associated symptoms
    \item \textbf{Examination:} the areas affected and signs of an underlying condition
    \item \textbf{Blood tests:} thyroid function, blood glucose, and other tests as indicated
    \item \textbf{Starch-iodine test:} can map the areas of sweating
    \item \textbf{Assessment of impact:} how the sweating affects daily life
\end{itemize}

\section*{Management}
Treatment is effective and stepped.
\begin{itemize}
    \item \textbf{Antiperspirants:} high-strength aluminium chloride antiperspirants applied at night
    \item \textbf{Iontophoresis:} a device that passes a mild current through the skin, effective for hands and feet
    \item \textbf{Botulinum toxin injections:} block the nerves to the sweat glands in the armpits and other areas, providing months of relief
    \item \textbf{Medicines:} anticholinergic tablets reduce sweating but have side effects
    \item \textbf{Laser and microwave treatments:} newer options that destroy sweat glands
    \item \textbf{Surgery:} rarely, removal of the sweat glands or interruption of the nerve supply
    \item \textbf{Treating the cause:} management of thyroid disease, menopause, or other underlying conditions
\end{itemize}

\section*{Complications}
\begin{itemize}
    \item Skin irritation and infections from constant wetness
    \item Embarrassment, social avoidance, and reduced self-esteem
    \item Effects on work and relationships
    \item Dehydration and heat intolerance in severe cases
    \item Anxiety about sweating, which worsens it
\end{itemize}

\section*{Prognosis and Outcomes}
Most people with hyperhidrosis improve significantly with treatment. Antiperspirants, iontophoresis, and botulinum toxin are effective, and the condition does not cause lasting harm. Treatment is tailored to the severity and the areas affected, and most people find an option that gives them good control.

\section*{Prevention and Self-Care}
\begin{itemize}
    \item Use a strong antiperspirant at night, not a deodorant
    \item Wear breathable fabrics such as cotton
    \item Shower daily and dry thoroughly, especially between the toes
    \item Avoid triggers such as spicy food, caffeine, and alcohol
    \item Manage stress with relaxation techniques
    \item Seek medical help rather than enduring the sweating
    \item Treat any skin infections promptly
\end{itemize}

\section*{Sources and Further Reading}
The following reputable sources provide further information for healthcare students and patients seeking reliable, current advice about excessive sweating.
\begin{itemize}
    \item DermNet. \textit{Hyperhidrosis.} https://dermnetnz.org/
    \item Healthdirect Australia. \textit{Excessive sweating (hyperhidrosis).} https://www.healthdirect.gov.au/
    \item International Hyperhidrosis Society. \textit{Information and treatment options.} https://www.sweathelp.org/
    \item NHS. \textit{Hyperhidrosis.} https://www.nhs.uk/conditions/hyperhidrosis/
    \item Mayo Clinic. \textit{Hyperhidrosis.} https://www.mayoclinic.org/diseases-conditions/hyperhidrosis/
    \item MSD Manual. \textit{Hyperhidrosis.} https://www.msdmanuals.com/
\end{itemize}
